% presentation/preamble.tex
\documentclass[aspectratio=169,11pt]{beamer}

\usetheme{metropolis}

\usepackage[french]{babel}
\usepackage[T1]{fontenc}
\usepackage{amsmath,amssymb}
\usepackage{graphicx}
\usepackage{xcolor}
\usepackage{booktabs} 
\usepackage{tikz} 

% ==============================================================================
% PALETTE DE COULEURS PREMIUM (Style Infographie Minimaliste)
% ==============================================================================
\definecolor{PremiumDark}{RGB}{45,40,35}     % Un brun-noir très profond pour le texte et titres (lisibilité 100%)
\definecolor{SubtleTaupe}{RGB}{115,100,85}   % Un taupe élégant pour la structure secondaire
\definecolor{Cream}{RGB}{253,251,247}        % Fond crème doux, haut de gamme, qui repose les yeux
\definecolor{LuxuryGold}{RGB}{180,145,90}    % Or brossé pour les alertes, barres et puces
\definecolor{LightCard}{RGB}{245,240,230}    % Couleur de fond pour les blocs/boîtes de texte

% ==============================================================================
% CONFIGURATION DES COULEURS DE BASE
% ==============================================================================
\setbeamercolor{background canvas}{bg=Cream}
\setbeamercolor{normal text}{fg=PremiumDark}
\setbeamercolor{title}{fg=PremiumDark}
\setbeamercolor{subtitle}{fg=SubtleTaupe}
\setbeamercolor{frametitle}{fg=PremiumDark}
\setbeamercolor{structure}{fg=PremiumDark}
\setbeamercolor{alerted text}{fg=LuxuryGold}

% ==============================================================================
% DESIGN DU PIED DE PAGE (Nettoyé et Ultra-Clair)
% ==============================================================================
\metroset{progressbar=foot} 
\setbeamercolor{progress bar}{fg=LuxuryGold, bg=LightCard} % Barre fine Or sur fond discret

% Le texte du bas écrit en gris-brun fin directement sur le fond Crème
\setbeamercolor{footline}{fg=SubtleTaupe, bg=Cream} 
\setbeamercolor{page number in head/foot}{fg=SubtleTaupe}

% ==============================================================================
% CORRECTION CRITIQUE : Suppression des pages sombres de Metropolis
% ==============================================================================
% On force Metropolis à faire des pages de garde et de sections CLAIRES. 
% Plus aucun risque de texte noir caché sur fond marron !
\setbeamercolor{palette primary}{fg=PremiumDark, bg=Cream}

% ==============================================================================
% INFOGRAPHIE : Structure visuelle des Diapositives
% ==============================================================================
\setbeamertemplate{navigation symbols}{}

% Ligne de séparation or brossé sous le titre de chaque slide
\setbeamertemplate{frametitle}{%
    \begin{beamercolorbox}[wd=\paperwidth,ht=4ex,dp=1.5ex,leftskip=0.5cm]{frametitle}
        \usebeamerfont{frametitle}\insertframetitle
    \end{beamercolorbox}%
    \vspace{-0.1cm}
    {\color{LuxuryGold}\hrule height 0.8pt} 
    \vspace{0.2cm}
}

% Blocs Modernes style "Cards / Encadrés épurés"
\setbeamertemplate{blocks}[rounded][shadow=false]
\setbeamercolor{block title}{bg=SubtleTaupe, fg=Cream}
\setbeamercolor{block body}{bg=LightCard, fg=PremiumDark}

\setbeamercolor{block title alerted}{bg=LuxuryGold, fg=Cream}
\setbeamercolor{block body alerted}{bg=LuxuryGold!10, fg=PremiumDark}

% Puces d'infographie épurées
\setbeamertemplate{itemize item}{\color{LuxuryGold}$\blacktriangleright$}
\setbeamertemplate{itemize subitem}{\color{SubtleTaupe}$\bullet$}

% Typographie Mathématique Professionnelle
\usefonttheme{professionalfonts}

\newcommand{\R}{\mathbb{R}}
\newcommand{\E}{\mathbb{E}}
\DeclareMathOperator*{\argmax}{arg\,max}